\chapter{绪论}
\label{chap:introduction}

\section{粒子物理与标准模型}
\label{sec:onepone}

物理是研究世界、理解世界、改变世界的学问。人类对世界运行规则的理解从牛顿力学，到分析力学，再到量子力学，以至量子场论，一路走过。量子场论作为物理发展史中的一枚皇冠，受到无数新生学者的追求和膜拜。粒子物理就是这样一门在量子场论基础上发展出来的学科。

粒子物理研究世界的本质。狭义地说，它研究构成万事万物的费米子、在其中传递相互作用的媒介子，以及赋予部分物质质量的上帝粒子。

标准模型是目前粒子物理学科中最广泛适用的理论框架，它的建立和不断完善汇聚了一代又一代物理学家的智慧和努力，其中包括盖尔曼、茨威格等人提出的夸克模型[1]，布罗特、恩格勒、希格斯等人提出的希格斯机制[2]，格拉肖、萨拉姆、温伯格等人提出的电弱统一理论[3]等。

具体来说，标准模型所描述的基本粒子有：三代夸克，包含上、下、奇、粲、顶、底夸克；三代轻子，包含电子、电子中微子、缪子、缪子中微子、陶子、陶子中微子；媒介子，包含光子、胶子、W及Z玻色子；以及希格斯粒子。此外还包含各基本粒子所对应的反粒子。标准模型所描述的基本相互作用有：电磁相互作用、强相互作用、弱相互作用。然而，引力相互作用尚未被纳入到该框架中，如何将其与标准模型统一仍是现代物理学未解决的难题。

标准模型的“最后一块拼图”希格斯粒子于2012年被欧洲大型强子对撞机上运行的两个探测器ATLAS[4]和CMS[5]所发现。回望历史，一个必须承认的事实是，物理学中辉煌成果的取得越来越依赖于全球物理学家的通力合作。

\section{量子纠缠与高能物理的结合}
\label{sec:oneptwo}

量子纠缠是量子力学中具有深刻物理意义的现象。在量子信息领域，纠缠是量子计算和量子通信的关键资源。

从本质上来说，量子力学提供了解释世界的全新视角以及可行的数学工具，因此将量子力学的观点应用到高能物理研究中是极为自然的。在过去，受限于实验装置和分析手段，在对撞机实验的高能标环境下对量子纠缠的研究尚未能够广泛地铺展开。如今，伴随着人工智能和机器学习浪潮的兴起以及蓬勃发展，越来越多针对海量数据进行处理的有效手段被开发出来。机器学习方法在高能物理中的应用也进一步催化了高能条件下对于量子纠缠现象的研究。2024年7月，大型强子对撞机上的ATLAS探测器首次发现顶夸克对的量子纠缠[6]，展现了利用对撞机探索量子力学基础问题的巨大潜力。

目前，高能对撞机环境下对量子纠缠现象的研究关注点集中在出射粒子的自旋上[7]。这主要是因为对于自旋，量子层析技术已为其提供了足够可行有效的计算范式[8]。利用衰变末态产物的运动学信息能够完整建立出射粒子对的自旋关联矩阵信息，进而使对纠缠存在与否的判断成为可能。

\section{研究主要内容和意义}
\label{sec:onepthree}

尽管在如上的交叉领域已开展了部分研究，但在大型强子对撞机上对于$ \tau^+ \tau^- $系统的纠缠的检验尚存在困难。这主要是因为$ \tau $子在衰变过程中会产生中微子，而中微子难以与对撞机中的探测器发生反应，导致中微子运动学信息的缺失，从而给$ \tau^+ \tau^- $系统自旋密度矩阵信息的重建造成障碍。

在该背景下，本研究主要关注在质子-质子对撞过程中产生的$ \tau^+ \tau^- $粒子之间的纠缠，利用机器学习手段重建量子层析所需要的末态粒子的运动学信息，最终计算该系统自旋关联矩阵的各矩阵元，尝试给出关于其纠缠存在与否的定性判断。该研究对量子力学基础问题的回答具有重要意义，在一定程度上能为后人的探索带来启发。

重建过程中所使用的机器学习方法主要包括扩散模型[9][10]和Transformer模型[11]，以及在二者基础上构建获得的OmniLearn生成模型[12][7]。机器学习与粒子物理的深度融合越来越成为分析手段创新的趋势，本研究所采用的架构相比传统的丢失质量计算器(Missing Mass Calculator)方法[13]具备明显优势，具有一定的推广意义。

\section{论文结构排布}
\label{sec:onepfour}

论文的剩余部分将按如下结构展开：第二章简述量子纠缠和量子层析方法的基本原理，以及在$ pp\rightarrow \tau^+ \tau^- $过程中的具体应用；第三章简述扩散模型和Transformer模型的基本原理，介绍OmniLearn生成模型的具体架构和相应的训练过程；第四章展示利用机器学习方法得到的各项结果，并与MMC方法进行比较；第五章是全文的总结，并重点阐述下一步工作计划。

